\chapter{Hướng dẫn tái hiện và bảng bổ sung}
\label{chap:phu-luc}

\section{Bốn mức tái hiện}
\label{sec:tai-hien}

\subsection{Từ dữ liệu đến bảng}

Toàn bộ mã nguồn, dữ liệu dẫn xuất và hướng dẫn nằm tại \url{https://github.com/doanan293/grounded-or-guessed}
(tệp \texttt{REPRODUCE.md}). Bốn mức, từ rẻ đến đắt:
\begin{enumerate}
\item \textbf{Mức 0} --- sinh lại bảng, hình và macro số liệu từ các tệp CSV đã có: \texttt{python scripts/make\_tables.py}
và \texttt{python scripts/make\_numbers.py}.
\item \textbf{Mức 1} --- sinh lại mọi tệp CSV từ gói kết quả các lượt chạy: giải nén gói rồi chạy
\texttt{scripts/regenerate\_all.sh}; trong Docker, các tệp CSV và bảng sinh ra trùng từng byte với bản đã công bố.
\item \textbf{Mức 2} --- chạy lại chính SemLoc và các phương pháp đối chứng, phát lại mọi câu trả lời LLM từ bộ đệm, không gọi mạng
(\texttt{FLX\_LLM\_OFFLINE=1}).
\item \textbf{Mức 3} --- chạy mới qua một điểm cuối LLM tương thích OpenAI, với mã lượt chạy mới.
\end{enumerate}

\subsection{Môi trường}

Ảnh Docker (\texttt{docker build -t grounded-or-guessed .}) chứa Python 3.12 và toàn bộ phụ thuộc, đặt chế độ không mạng
cho LLM; lệnh mặc định chạy bộ kiểm thử đơn vị. Gói kết quả các lượt chạy (đầu ra SemLoc, nhật ký và bộ đệm lượt gọi LLM,
bảng băm kiểm tra) được tạo bằng \texttt{scripts/export\_artifact.py}; công cụ từ chối đóng gói nếu phát hiện khóa truy
cập hay địa chỉ thư điện tử.

\section{Bảng bổ sung}
\label{sec:bang-bo-sung}

\subsection{Phân rã trên nhật ký công bố}

\begin{table}[H]
\centering
\caption{Phân rã Shapley trên nhật ký công bố của tác giả (dòng lỗi chuẩn gốc)}
\label{tab:phan-ra-cong-bo}
\resizebox{\linewidth}{!}{\begin{tabular}{llrrrr}
\toprule
Cấu hình & Tập con & Chỉ số & SemLoc & Điểm thưởng dòng đoán & Phổ ngữ nghĩa \\
\midrule
SemLoc, Claude Sonnet 4,6, $T{=}0,8$ & tất cả & Acc@1 & 42,8 & 27,8{\scriptsize\,[22,5; 33,5]} & 15,0{\scriptsize\,[12,1; 18,0]} \\
SemLoc, Claude Sonnet 4,6, $T{=}0,8$ & tất cả & Acc@3 & 65,6 & 26,4{\scriptsize\,[21,9; 31,4]} & 39,2{\scriptsize\,[34,2; 44,1]} \\
SemLoc, Claude Sonnet 4,6, $T{=}0,8$ & thực thi được & Acc@1 & 57,8 & 37,6{\scriptsize\,[31,0; 44,3]} & 20,3{\scriptsize\,[16,3; 24,3]} \\
SemLoc, Claude Sonnet 4,6, $T{=}0,8$ & thực thi được & Acc@3 & 77,3 & 35,7{\scriptsize\,[30,0; 41,6]} & 41,6{\scriptsize\,[36,3; 47,2]} \\
SemLoc, Gemini 2,5 Pro, $T{=}0,8$ & tất cả & Acc@1 & 39,6 & 26,0{\scriptsize\,[21,1; 31,3]} & 13,6{\scriptsize\,[10,8; 16,6]} \\
SemLoc, Gemini 2,5 Pro, $T{=}0,8$ & tất cả & Acc@3 & 64,4 & 24,2{\scriptsize\,[20,0; 28,9]} & 40,2{\scriptsize\,[34,7; 45,4]} \\
SemLoc, Gemini 2,5 Pro, $T{=}0,8$ & thực thi được & Acc@1 & 53,5 & 35,1{\scriptsize\,[28,6; 42,2]} & 18,4{\scriptsize\,[14,7; 22,1]} \\
SemLoc, Gemini 2,5 Pro, $T{=}0,8$ & thực thi được & Acc@3 & 73,0 & 32,7{\scriptsize\,[27,3; 38,5]} & 40,3{\scriptsize\,[34,8; 45,6]} \\
\bottomrule
\end{tabular}
}
\end{table}

\subsection{Kết quả đầy đủ trên ba bộ nhãn}

\begin{table}[H]
\centering
\caption{RQ2 đầy đủ: mọi biến thể trên ba bộ nhãn (tập test, Gemini 3.8 Flash, 3 lần chạy; $\Delta$ so với SemLoc)}
\label{tab:rq2-day-du}
\adjustbox{max width=\linewidth, max totalheight=\TallTableHeight}{\begin{tabular}{lrrrrrr}
\toprule
Phương pháp & Acc@1 & Acc@3 & Acc@5 & $\Delta$Acc@1 [CI 95\%] & $\Delta$Acc@3 [CI 95\%] & $r_{\text{rb}}$ \\
\midrule
\multicolumn{7}{l}{\emph{Ground truth gốc, $n=200$, 3 lần lặp}} \\
SemLoc (bản phát hành) & 36,8{\scriptsize$\pm$0,8} & 60,3{\scriptsize$\pm$2,3} & 75,0{\scriptsize$\pm$0,9} & -- & -- & -- \\
\quad bỏ điểm thưởng dòng đoán & 29,2{\scriptsize$\pm$1,4} & 57,7{\scriptsize$\pm$2,5} & 72,7{\scriptsize$\pm$1,3} & -7,7{\scriptsize\,[-13,9; -1,6]} & -2,7{\scriptsize\,[-4,9; -0,9]} & -0,38 \\
\quad chỉ thưởng khi có căn cứ & 34,7{\scriptsize$\pm$1,5} & 58,2{\scriptsize$\pm$2,9} & 72,8{\scriptsize$\pm$1,5} & -2,2{\scriptsize\,[-4,2; -0,6]} & -2,2{\scriptsize\,[-4,2; -0,6]} & -0,94 \\
\quad + chỉ cổng độ phủ & 36,8{\scriptsize$\pm$0,8} & 57,7{\scriptsize$\pm$1,8} & 62,5{\scriptsize$\pm$1,3} & +0,0{\scriptsize\,[0,0; 0,0]} & -2,7{\scriptsize\,[-7,7; 1,7]} & -0,59 \\
\quad + chỉ $\lambda$·Ochiai & 36,8{\scriptsize$\pm$0,8} & 61,3{\scriptsize$\pm$1,9} & 75,0{\scriptsize$\pm$0,9} & +0,0{\scriptsize\,[0,0; 0,0]} & +1,0{\scriptsize\,[0,2; 2,3]} & +0,54 \\
\quad + xếp lại bằng CF & 35,0{\scriptsize$\pm$1,0} & 62,2{\scriptsize$\pm$2,5} & 77,7{\scriptsize$\pm$1,8} & -1,8{\scriptsize\,[-3,8; -0,3]} & +1,8{\scriptsize\,[0,5; 3,5]} & +0,36 \\
CovGate (nhóm đề xuất) & 36,8{\scriptsize$\pm$0,8} & 61,3{\scriptsize$\pm$1,9} & 75,0{\scriptsize$\pm$0,9} & +0,0{\scriptsize\,[0,0; 0,0]} & +1,0{\scriptsize\,[0,2; 2,3]} & +0,54 \\
CovGate + xếp lại bằng CF & 35,0{\scriptsize$\pm$1,0} & 62,5{\scriptsize$\pm$2,2} & 77,7{\scriptsize$\pm$1,8} & -1,8{\scriptsize\,[-3,8; -0,3]} & +2,2{\scriptsize\,[0,9; 3,8]} & +0,46 \\
LLM hỏi thẳng (1 prompt) & 40,0{\scriptsize$\pm$0,0} & 43,0{\scriptsize$\pm$0,5} & 43,2{\scriptsize$\pm$0,8} & +3,2{\scriptsize\,[-0,8; 7,4]} & -17,3{\scriptsize\,[-23,3; -11,1]} & -0,51 \\
LLM + độ phủ (1 prompt, nhóm đề xuất) & 40,2{\scriptsize$\pm$0,6} & 68,2{\scriptsize$\pm$0,3} & 70,0{\scriptsize$\pm$0,9} & +3,3{\scriptsize\,[0,5; 6,5]} & +7,8{\scriptsize\,[0,2; 15,2]} & +0,04 \\
SBFL Ochiai & 3,0{\scriptsize$\pm$0,0} & 15,5{\scriptsize$\pm$0,0} & 48,0{\scriptsize$\pm$0,0} & -33,8{\scriptsize\,[-41,2; -26,8]} & -44,8{\scriptsize\,[-52,8; -36,6]} & -0,87 \\
\midrule
\multicolumn{7}{l}{\emph{Ground truth đã sửa bằng thực thi, $n=179$, 3 lần lặp}} \\
SemLoc (bản phát hành) & 52,3{\scriptsize$\pm$0,9} & 75,4{\scriptsize$\pm$2,0} & 82,7{\scriptsize$\pm$1,1} & -- & -- & -- \\
\quad bỏ điểm thưởng dòng đoán & 43,2{\scriptsize$\pm$2,1} & 70,8{\scriptsize$\pm$2,3} & 78,4{\scriptsize$\pm$1,7} & -9,1{\scriptsize\,[-16,7; -1,9]} & -4,7{\scriptsize\,[-7,9; -2,1]} & -0,37 \\
\quad chỉ thưởng khi có căn cứ & 48,2{\scriptsize$\pm$1,7} & 71,3{\scriptsize$\pm$2,8} & 78,6{\scriptsize$\pm$2,0} & -4,1{\scriptsize\,[-7,2; -1,7]} & -4,1{\scriptsize\,[-7,2; -1,7]} & -1,00 \\
\quad + chỉ cổng độ phủ & 52,3{\scriptsize$\pm$0,9} & 79,1{\scriptsize$\pm$2,0} & 85,1{\scriptsize$\pm$1,3} & +0,0{\scriptsize\,[0,0; 0,0]} & +3,7{\scriptsize\,[1,5; 6,6]} & +1,00 \\
\quad + chỉ $\lambda$·Ochiai & 52,3{\scriptsize$\pm$0,9} & 79,3{\scriptsize$\pm$1,7} & 85,1{\scriptsize$\pm$1,3} & +0,0{\scriptsize\,[0,0; 0,0]} & +3,9{\scriptsize\,[1,6; 6,8]} & +0,85 \\
\quad + xếp lại bằng CF & 49,9{\scriptsize$\pm$0,9} & 76,2{\scriptsize$\pm$1,8} & 83,4{\scriptsize$\pm$1,4} & -2,4{\scriptsize\,[-4,8; -0,6]} & +0,7{\scriptsize\,[0,0; 2,1]} & -0,35 \\
CovGate (nhóm đề xuất) & 52,3{\scriptsize$\pm$0,9} & 79,3{\scriptsize$\pm$1,7} & 85,1{\scriptsize$\pm$1,3} & +0,0{\scriptsize\,[0,0; 0,0]} & +3,9{\scriptsize\,[1,6; 6,8]} & +0,85 \\
CovGate + xếp lại bằng CF & 49,9{\scriptsize$\pm$0,9} & 77,1{\scriptsize$\pm$2,2} & 85,3{\scriptsize$\pm$1,2} & -2,4{\scriptsize\,[-4,8; -0,6]} & +1,7{\scriptsize\,[0,5; 3,3]} & +0,28 \\
LLM hỏi thẳng (1 prompt) & 56,8{\scriptsize$\pm$0,7} & 62,8{\scriptsize$\pm$0,7} & 62,9{\scriptsize$\pm$0,9} & +4,5{\scriptsize\,[-0,4; 9,7]} & -12,7{\scriptsize\,[-19,7; -6,2]} & -0,28 \\
LLM + độ phủ (1 prompt, nhóm đề xuất) & 57,2{\scriptsize$\pm$0,9} & 92,0{\scriptsize$\pm$0,3} & 93,7{\scriptsize$\pm$0,9} & +4,8{\scriptsize\,[1,5; 8,7]} & +16,6{\scriptsize\,[10,5; 22,8]} & +0,65 \\
SBFL Ochiai & 3,9{\scriptsize$\pm$0,0} & 19,0{\scriptsize$\pm$0,0} & 67,6{\scriptsize$\pm$0,0} & -48,4{\scriptsize\,[-58,6; -38,3]} & -56,4{\scriptsize\,[-67,0; -45,4]} & -0,85 \\
\bottomrule
\end{tabular}
}
\end{table}

\section{Các prompt}
\label{sec:prompt}

\subsection{LLM + kiểm thử thất bại (phương pháp đề xuất) và biến thể có độ phủ}

Thông điệp hệ thống yêu cầu trả lời đúng định dạng JSON. Biến thể có độ phủ thêm dòng giải thích chú thích (đánh dấu
\texttt{[cov]} dưới đây) và chú thích \texttt{[f=…/p=…]} ở cuối mỗi dòng mã; mọi phần khác giống hệt.
\begin{lstlisting}[numbers=none]
You are a fault-localization assistant. Reply with JSON only:
{"lines": [<line numbers, most suspicious first, at most 5>]}

The program below has exactly one faulty line. <n_fail> tests fail and <n_pass> pass.
[cov] Each line is annotated with [f=<number of failing tests that execute it>/p=<number of passing tests>].
<numbered source>
Failing tests:
- <test name>: <first error line>
Return the five most suspicious line numbers, most suspicious first.
\end{lstlisting}

\subsection{Đề xuất bản sửa (bộ phán định nhãn và đề xuất rồi kiểm chứng)}

\begin{lstlisting}[numbers=none]
You fix one-line bugs in Python. Reply with JSON only:
{"kind": "replace" | "insert_after", "code": "<one line of Python, no indentation>"}

Program (line-numbered): <numbered source>
Failing tests: <pytest failure summary>
Assume the bug is at line <n>. Give the single line that should replace line <n> (or, if code is missing,
the line to insert after line <n>) so that all tests pass. Make the smallest possible change (ideally one
or two tokens) and keep everything else identical. Attempt <k>.
\end{lstlisting}

\section{Quy trình duyệt nhãn}
\label{sec:quy-trinh-duyet}

Mỗi thành viên nhận một tệp gồm các ca được giao (\texttt{docs/analysis/gt-review/duyet-lan-1-co-goi-y/}). Mỗi ca có liên kết tới
chương trình và kiểm thử, đoạn mã quanh các dòng ứng viên, và nhận định của LLM để tham khảo; các dòng ứng viên được ký
hiệu G (nhãn gốc), S (dòng bộ phán định dời tới) và F (dòng LLM đề xuất). Người duyệt ghi vào cột \texttt{answer} một hay
nhiều chữ (ví dụ \texttt{G+F} cho lỗi hai dòng) hoặc một số dòng khác, và có thể ghi \texttt{?} nếu không chắc. Bộ nhãn
đã rà soát được dựng tự động từ các tệp này (\texttt{scripts/build\_reviewed\_gt.py}); các ca trưởng nhóm sửa được liệt
kê trong mã kèm lý do.
